\chapter{哈希表结构与算法设计}
\label{ch:system}

本章依照系统核心设定，给出分层哈希表各模块组件的完整结构与算法细节。内容覆盖 Facility 、 MiniArray、多层级 Cubby、Cubby 内 k-kick 与 Local Query Router，商压缩的实现，并在最后一节统一说明插入、查询与删除的操作语义。

\section{Facility 与 MiniArray}

\subsection{Facility 的基本属性}

每个 Facility 对应全局表中的一个局部片段，规模为 \(K=\mathrm{polylog}\,N\)。Facility 管理其内部的全部 Cubby（包括各 tier 及唯一 tail），并维护 一个 Facility 层次的 MiniArray，用于定位键对应的 Cubby 。B 树的 fanout(上文有定义，fanout 为每个内部节点的子指针个数) 记为 \(F\)，取 \(F=\Theta(\log n/\log\log n)\)；树高为 \(O(1)\)，每个内部节点的子指针仅需 \(\Theta(\log\log n)\) 比特\cite{bender2022otsh,bayer1972,pagh2001}。

\subsection{MiniArray 的逻辑结构}

MiniArray 由以下部分共同组成：
\begin{center}
  \small
  MiniArray $\rightarrow$ $O(1)$ 深度 B 树 $\rightarrow$ Packed Leaf + Bitmap + Rank/Select + Lookup Table
\end{center}

叶节点 Leaf 包含：
\begin{itemize}
  \item bitmap：用于标记叶内哪些逻辑位置已有有效元素，其大小与叶容量相同；
  \item data：PackedArray，在有效位置上紧凑、无间隙地存放变长元素。
\end{itemize}

内部节点 Internal 包含：
\begin{itemize}
  \item subtree\_size$[F]$：每个子树中有效元素总数，供 Rank/Select 定位；
  \item children$[F]$：至多 \(F\) 个子女指针，每个指针 \(\Theta(\log\log n)\) 比特。
\end{itemize}

MiniArray 对外提供的接口函数为 Access(idx)、Insert(idx, val)、Delete(idx)、Update(idx, val)，其中 idx 为逻辑下标。叶容量上限为 MAX\_LEAF\_SIZE，由 NODE\_MAX\_BITS 常数 \(c\in[1,3]\) 产生\cite{bender2022otsh,kurpicz2023pachash,lehmann2024shockhash}。

\subsection{Access：按逻辑下标读取}

Access 操作的输入是逻辑下标，搜索过程自根向下逐层向下遍历，利用各层 subtree\_size 进行前缀和判断，将输入的逻辑下标映射到对应的目标叶节点；在叶子节点的内部用 Select 在 bitmap 上定位第 idx 个有效位置，再从 data 读出元素。树高为 \(O(1)\)，叶内 Select 算法通过 Lookup Table 常数时间完成\cite{jacobson1989,pagh2008retrieval,kurpicz2023pachash}， 不需要遍历比特串，获取到物理位置之后，直接从紧凑数组中读取对应的元素，并返回。

\begin{algorithm}[htbp]
  \caption{MiniArray.Access}
  \label{alg:miniarray-access}
  \begin{algorithmic}[1]
    \State 输入逻辑下标 idx
    \State cur $\gets$ root，rem $\gets$ idx
    \While{cur 为内部节点}
      \State 找到子树 $i$ 使 $\displaystyle\sum_{t<i}\mathit{subtree\_size}[t]\le \mathit{rem}<\sum_{t\le i}\mathit{subtree\_size}[t]$
      \State rem $\gets$ rem - $\displaystyle\sum_{t<i}\mathit{subtree\_size}[t]$
      \State cur $\gets$ cur.child[i]
    \EndWhile
    \State pos $\gets$ Call{Select}(cur.bitmap, rem)
    \State 返回 cur.data[pos]
  \end{algorithmic}
\end{algorithm}

\subsection{Insert：插入与叶分裂}

插入过程中，Insert 在定位叶节点的过程中维护路径栈，为后续自底向上更新 subtree\_size 提供途径；在叶节点内用 Select 函数找到插入位置，将 bitmap 对应位置赋值为 1，并向 data 写入 val。若叶内有效元素数达到叶子容量的上限，就会调用 Split 对叶子节点进行分裂操作；插入与分裂检查完成之后，自底向上对路径栈中每个内部节点的 subtree\_size 加一， 以保证后续查询与定位的正确性。

\begin{algorithm}[htbp]
  \caption{MiniArray.Insert}
  \label{alg:miniarray-insert}
  \begin{algorithmic}[1]
    \State 输入 idx, val；初始化空路径栈 stk
    \State cur $\gets$ root，rem $\gets$ idx
    \While{cur 为内部节点}
      \State stk.push(cur)
      \State 找到子树 $i$ 使 $\displaystyle\sum_{t<i}\mathit{subtree\_size}[t]\le \mathit{rem}<\sum_{t\le i}\mathit{subtree\_size}[t]$
      \State rem $\gets$ rem - $\displaystyle\sum_{t<i}\mathit{subtree\_size}[t]$
      \State cur $\gets$ cur.child[i]
    \EndWhile
    \State leaf $\gets$ cur
    \State pos $\gets$ Call{Select}(leaf.bitmap, rem)
    \State leaf.bitmap.set(pos)；leaf.data.insert(pos, val)
    \If{leaf.size $\ge$ MAX\_LEAF\_SIZE}
      \State Call{Split}(leaf)
    \EndIf
    \For{路径栈中每个内部节点 node}
      \State node.subtree\_size[i] $\mathrel{+}=1$
    \EndFor
  \end{algorithmic}
\end{algorithm}

\subsection{Delete：删除与叶合并}

Delete 同样沿 subtree\_size 定位叶节点，用 Select 函数在 Bitmap 上找到待删元素的具体位置，清除 bitmap 对应位，并从 data 移除对应的元素。若叶内有效元素过少，低于阈值，调用 Merge 与相邻兄弟叶合并，合并可能需要多次操作才能保证 B 树恢复数据结构；最后自底向上将路径上 subtree\_size 减一。

\begin{algorithm}[htbp]
  \caption{MiniArray.Delete}
  \label{alg:miniarray-delete}
  \begin{algorithmic}[1]
    \State 输入 idx；初始化路径栈 stk
    \State 沿 subtree\_size 定位至叶 leaf(同 Insert)
    \State pos $\gets$ Call{Select}(leaf.bitmap, rem)
    \State leaf.bitmap.unset(pos)；leaf.data.remove(pos)
    \If{leaf 有效元素过少（即 leaf.size < MAX\_LEAF\_SIZE / 2）}
      \State Call{Merge}(leaf)
    \EndIf
    \For{路径栈中每个内部节点 node}
      \State node.subtree\_size[i] $\mathrel{-}=1$
    \EndFor
  \end{algorithmic}
\end{algorithm}

\subsection{Update：原地覆写}

Update 仅修改已有有效位置的元素值，不改变 bitmap 中 1 的个数，也不触发 Split/Merge。定位过程与 Access 相同，在叶内 Select 后直接 leaf.data[pos] $\gets$ val。

\begin{algorithm}[htbp]
  \caption{MiniArray.Update}
  \label{alg:miniarray-update}
  \begin{algorithmic}[1]
    \State 输入 idx, val
    \State 沿 subtree\_size 定位至叶 leaf
    \State pos $\gets$ Call{Select}(leaf.bitmap, rem)
    \State leaf.data[pos] $\gets$ val
  \end{algorithmic}
\end{algorithm}

\subsection{Split 与 Merge}

Split（叶分裂）。 新建叶节点，以 leaf.size/2 为界将后半段 bitmap 与 data 迁至新叶，原叶保留前半段；在父节点 child 数组中插入新叶指针，并更新 subtree\_size 的分裂计数。

\begin{algorithm}[htbp]
  \caption{MiniArray.Split（叶分裂）}
  \label{alg:miniarray-split}
  \begin{algorithmic}[1]
    \State 输入叶 leaf
    \State new\_leaf $\gets$新建 LeafNode；mid $\gets$ leaf.size/2
    \State 将 leaf.bitmap、leaf.data 后半段搬至 new\_leaf
    \State leaf 保留前半段
    \State 在父节点 parent.child 中插入 new\_leaf，更新 parent.subtree\_size
  \end{algorithmic}
\end{algorithm}

Merge（叶合并）。 取相邻兄弟叶，将兄弟的 bitmap 与 data 并入当前叶，从父节点删去兄弟子树并刷新 subtree\_size。

\begin{algorithm}[htbp]
  \caption{MiniArray.Merge（叶合并）}
  \label{alg:miniarray-merge}
  \begin{algorithmic}[1]
    \State 输入叶 leaf；sib $\gets$相邻兄弟叶；parent $\gets$ leaf 的父节点
    \State leaf.bitmap $\mathrel{|=}$ sib.bitmap；leaf.data.append(sib.data)
    \State 从 parent.child 移除 sib，更新 parent.subtree\_size
  \end{algorithmic}
\end{algorithm}

\subsection{Rank 与 Select}

叶子节点中 bitmap 上的 Rank/Select 通过 Lookup Table 实现，这是紧凑动态结构的标准组件\cite{jacobson1989,pagh2008retrieval,bender2022otsh,lehmann2026mphsurvey}：
\begin{itemize}
  \item Rank(bitmap, \(i\))：返回 bitmap 前 \(i\) 位中 1 的个数，等于 LookupTable 对高段与低段的查表结果之和；
  \item Select(bitmap, \(k\))：返回 bitmap 中第 \(k\) 个 1 的位置，由 LookupTable\([k]\) 直接给出。
\end{itemize}

\begin{algorithm}[htbp]
  \caption{Rank / Select（查表实现）}
  \label{alg:rank-select}
  \begin{algorithmic}[1]
    \Function{Rank}{bmp, $i$}
      \State 返回 LookupTable[bmp 高段] $+$ LookupTable[bmp 低段]
    \EndFunction
    \Function{Select}{bmp, $k$}
      \State 返回 LookupTable[$k$] \Comment{第 $k$ 个 1 的位置}
    \EndFunction
  \end{algorithmic}
\end{algorithm}

\subsection{复杂度与空间性质}

Facility 级 MiniArray 的树高为 \(O(1)\)，自根至叶的每层路由仅依赖 subtree\_size 比较，叶内 Rank/Select 经 Lookup Table 常数时间完成\cite{jacobson1989,pagh2008retrieval}。因此 Access 与 Update 为最坏情况 \(O(1)\)；Insert 与 Delete 在叶分裂/合并均摊意义下亦为 \(O(1)\)\cite{bayer1972,bender2022otsh}。空间方面，每个内部节点的 \(F\) 个子指针各需 \(\Theta(\log\log n)\) 比特，叶节点以 bitmap 标记有效位置、以 PackedArray 无间隙存放变长元素，整体辅助位数与 \(O(\log^{(k)}n)\) 的全局预算相容\cite{kurpicz2023pachash,lehmann2024shockhash}。Cubby 内的 \(A\)、\(M\)、\(B\) 三组 MiniArray 复用同一套接口语义：\(A[i]\) 存 router 映射，\(M\) 存 bin 空闲图，\(B\) 与 storage 槽位一一对应存差分编码——三者在 k-kick 搬移时须原子同步更新。

\section{多层级 Cubby 设计}

\subsection{层级属性与初始状态}

Facility 内 Cubby 按 tier 分级管理。\(j\)-tiered Cubby 的目标数量为
\[
  t_j=\Theta\!\left(\frac{(\log^{(j+1)}n)^2}{(\log^{(j)}n)^2}\right),
\]
容量（槽位数）为
\[
  r_j=\frac{K}{(\log^{(j)}n)^2}.
\]
系统初始时，每个 Facility 仅包含一个 1-tiered Cubby，且该 Cubby 即为 tail。此后所有新键插入 tail；tail 满后，旧 tail 被视为一个普通的 1-tiered Cubby 挂入 tier 集合，并新建 tail 承接后续插入，该过程可能触发 tier 向上合并\cite{bender2022otsh,bender2024sosa,bender2024iceberg,pandey2023iceberght}。

\subsection{动态调整：合并与拆分}

设当前 \(j\)-tiered Cubby 的实际个数为 \(\mathrm{cnt}_j\)。调度规则如下：
\begin{itemize}
  \item 当 \(\mathrm{cnt}_j\ge 3t_j\) 时，取出 \(t_j\) 个 \(j\)-tiered Cubby，将其内全部元素导出，合并为一个 \((j+1)\)-tiered Cubby 并重新插入；
  \item 当 \(\mathrm{cnt}_j\le t_j/2\) 且存在非空的 \((j+1)\)-tiered Cubby 时，取出一个 \((j+1)\)-tiered Cubby，拆成 \(t_j\) 个 \(j\)-tiered Cubby 并重插元素；
  \item 无论合并或拆分，均须取出 cubby 内所有元素重新插入，并同步更新对应 Cubby 中的 MiniArray \(A\)、\(M\)、\(B\) 与 storage，不能仅复制旧槽位内容（因 \(r_j\)、\(|I|\) 与商压缩上下文可能已变）\cite{bender2022otsh,conway2018,kuszmaul2024filters}。
\end{itemize}

\begin{algorithm}[htbp]
  \caption{j-tiered cubby 数量调度}
  \label{alg:tier-schedule}
  \begin{algorithmic}[1]
    \For{每个 tier 层级 \(j\)}
      \State 计算目标数量 \(t_j\)，统计实际数量 \(\mathrm{cnt}_j\)
      \If{\(\mathrm{cnt}_j\ge 3t_j\)}
        \State 取出 \(t_j\) 个 \(j\)-tiered Cubby，合并为一个 \((j+1)\)-tiered Cubby 并重新插入
      \ElsIf{\(\mathrm{cnt}_j\le t_j/2\) 且 tier \(j+1\) 非空}
        \State 取出一个 \((j+1)\)-tiered Cubby，拆成 \(t_j\) 个 \(j\)-tiered Cubby 并重新插入
      \EndIf
    \EndFor
  \end{algorithmic}
\end{algorithm}

\subsection{Cubby 内部组件}

每个容量为 \(|I|\) 的 Cubby 包含以下五部分：

storage 数组。 长度为 \(|I|\)，存放键值的商压缩载荷；槽位 \(i\) 上的 storage\([i]\) 对应 \(\pi(x)[\log N+1,\log U]\) 的 payload 部分，完整键恢复依赖 \(B[i]\) 与全局上下文。

k-kick 插入逻辑。 在 storage 上按一定的规则动态决定元素实际槽位，并维护每个占用元素的 insert\_depth。

MiniArray \(A\)。 长度为 \(|I|\) 的 Local Query Router 数组；\(A[i]\) 是一棵二叉前缀树，管理所有满足 \(g_I(x)=i\) 的键。

MiniArray \(M\)。 存储各 k-kick 深度 bin 的空闲槽位图，使插入、删除与 ReInsert 能在常数时间内判定 bin 是否含空位、定位空位，而不必扫描整个 Cubby。

MiniArray \(B\)。 与 storage 槽位一一对应。当 storage\([i]\) 存放键 \(x\) 的信息时，\(B[i]\) 保存
\[
  \delta=g_I(x)-i,\qquad g_I(x)=g^\ast(x)[\log|I|+1,\log K],
\]
以及 \(g^\ast(x)[\log|I|+1,\log K]\) 的中间位。配合 Facility 编号 \(r=g^\ast(x)[\log K+1,\log N]\) 与 payload，可完整恢复 \(\pi(x)\)\cite{bender2022otsh,pagh2008retrieval}。

\subsection{插入与删除操作}

插入。 每次插入均进入当前 Facility 的 tail Cubby。若 tail 仍有空位，则在 tail 内执行 k-kick 插入；若 tail 已满，则新建 tail，并将已满的旧 tail 视作 1-tiered Cubby 挂入 1-tier Cubby 集合，该步骤可能因 \(\mathrm{cnt}_1\ge 3t_1\) 而触发向上合并。

删除与 tail 回填。 若删除发生在 tail 之外的其他 Cubby，删除完成后从 tail 中\emph{随机}选取一个占用元素，经 k-kick 规则回填至刚空出的槽位，并更新 MiniArray \(A\)、\(M\)、\(B\) 与 storage。若 tail 在回填后为空，则从 1-tiered 层取出一个 Cubby 提升为 tail；该过程可能因 \(\mathrm{cnt}_i\le t_i\) 而触发向下拆分(甚至可能会触发多层)。上述不变量保证：除 tail 外，各 Cubby 尽量维持满载，局部波动由 tail 与 tier 调度共同吸收\cite{bender2022otsh,aamand2021dynamic,bansal2022deletions}。

\begin{algorithm}[htbp]
  \caption{tail 维护与删除回填}
  \label{alg:tail-ops}
  \begin{algorithmic}[1]
    \State 插入：写入 tail；若 tail 满则旧 tail 归入 1-tiered Cubby 并新建 tail
    \State 删除（非 tail）：释放被删槽位；从 tail 随机取元素 \(y\) 回填
    \State 对 \(y\) 执行 k-kick 写入，更新 \(A,M,B\)
    \If{tail 为空}
      \State 从 1-tiered Cubby 提升一个 Cubby 为 tail；必要时触发 tier 拆分
    \EndIf
  \end{algorithmic}
\end{algorithm}

\section{k-kick 层级插入算法}

约定：Cubby 内算法中的 \(x\) 均指 \(\pi(x)\)（置换后的值），而非原始键；符号仍记为 \(x\)，设 Cubby 容量为 \(|I|\)。

\subsection{层级 bin 定义}

k-kick 不对 storage 做物理分层，仅在逻辑上将 Cubby 划分为 \(k+1\) 个深度，\(k\) 为常数（通常 \(k\in[3,5]\)，基准取 4），定义不同层次的 bin：
\begin{itemize}
  \item 深度 0：唯一 bin 覆盖整个 Cubby，大小 \(s_0=K=\mathrm{polylog}\,n\)；
  \item 深度 1：将深度 0 的 bin 均匀切分为许多大小 \(s_1=\Theta((\log K)^6)\) 的子 bin；
  \item \( \dots \)
  \item 深度 \(i\)（\(i\ge 1\)）：bin 大小 \(s_i=\Theta((\log^{(i)}K)^6)\)，由上一层 bin 继续均匀切分得到。
\end{itemize}

定义函数 Bin(\(x,i\)) 返回键 \(x\) 在深度 \(i\) 所属 bin 的槽位区间：深度 0 时返回 \([0,|I|-1]\)；深度 \(i\ge 1\) 时在 Bin(\(x,i-1\)) 内按 \(s_i\) 均匀切分，并根据 \(x\) 的哈希比特在每一层中选定对应的一个子区间\cite{bender2022otsh,carter1979}。

\subsection{偏好序列}

对键 \(x\) 定义偏好序列
\[
  g_0(x)\rightarrow g_1(x)\rightarrow\cdots\rightarrow g_k(x).
\]
其中 \(g_0(x)=\operatorname{Bin}(x,0)\) 为整个 Cubby；\(g_{i+1}(x)\) 为 \(g_i(x)\) 经均匀切分后 \(x\) 所落子 bin。于是 \(g_{i+1}(x)\) 一定是 \(g_i(x)\) 的子 bin；沿路径自顶向下，每一层均为上一层的随机子区间。自底向上回溯 \(g_k,g_{k-1},\ldots,g_0\) 即可得到全部祖先 bin。

\begin{algorithm}[htbp]
  \caption{PreferenceSequence：偏好序列}
  \label{alg:pref-seq}
  \begin{algorithmic}[1]
    \Function{PreferenceSequence}{$x$}
      \State $g[0]\gets\Call{Bin}{x,0}$
      \For{$i=1$ to $k$}
        \State $g[i]\gets g[i-1]$ 内均匀切分后、$x$ 所属的子 bin
      \EndFor
      \State 返回 $g[0..k]$
    \EndFunction
  \end{algorithmic}
\end{algorithm}

\subsection{探测位置序列}

\(h_j(x)\) 表示 \(x\) 在 Cubby storage 中的第 \(j\) 个候选槽位。序列按先深后浅顺序构造：依次遍历 \(g_k(x),g_{k-1}(x),\ldots,g_0(x)\)，并在每个 bin 内按槽位顺序枚举全部位置。

对深度 \(i\) 的 bin \(g_i(x)\)，设
\[
  \mathrm{base}_i=s_k+s_{k-1}+\cdots+s_{i+1},
\]
则对任意 \(j\in[0,s_i-1]\)，全局探测序列中对应项为
\[
  h_{\mathrm{base}_i+j}(x)=g_i(x)\text{ 的第 }(j+1)\text{ 个槽位}.
\]
注意文档中下标公式 \(h_{(k+1-i)\cdot s_i+j}(x)\) 与 \(\mathrm{base}_i\) 表述等价：更深层的 bin 在序列中排在更前。序列总长度约为 \(\displaystyle\sum_{i=0}^{k}s_i\)；查询阶段并不线性扫描该序列，而由 Local Query Router 直接给出实际使用的下标 \(j\)\cite{bender2022otsh,pagh2004linear}。

\begin{algorithm}[htbp]
  \caption{ProbeSequence：探测位置序列（先深后浅）}
  \label{alg:probe-seq}
  \begin{algorithmic}[1]
    \Function{ProbeSequence}{$x$}
      \State $g\gets\Call{PreferenceSequence}{x}$；$\mathit{seq}\gets[\,]$
      \For{$i=k$ downto $0$}
        \State $\mathit{bin}\gets g[i]$
        \For{$\mathit{pos}$ 从 $\mathit{bin.start}$ 到 $\mathit{bin.end}$}
          \State $\mathit{seq.append(pos)}$
        \EndFor
      \EndFor
      \State 返回 $\mathit{seq}$ \Comment{$h_1(x),h_2(x),\ldots$ 依次为 $\mathit{seq}[0],\mathit{seq}[1],\ldots$}
    \EndFunction
  \end{algorithmic}
\end{algorithm}

\subsection{bin 饱和与最大可用深度}

在深度 \(d\) 下，称 bin 饱和，当且仅当 bin 已满，且其中每个元素 \(y\) 均满足 y.insert\_depth \(\ge\) d。bin 未满时不饱和；bin 已满但存在 y.insert\_depth < d 的元素时，亦不视为饱和。

GetMaxUsableDepth(\(x\)) 自 \(k\) 向 0 扫描偏好序列，返回最深的不饱和 bin 的深度；若全部饱和则返回 0(说明已满)。

\begin{algorithm}[htbp]
  \caption{IsSaturated 与 GetMaxUsableDepth}
  \label{alg:saturated}
  \begin{algorithmic}[1]
    \Function{IsSaturated}{bin, $d$}
      \If{bin 未满} \State 返回 false \EndIf
      \For{bin 中每个元素 $y$}
        \If{y.insert\_depth < d} \State 返回 false \EndIf
      \EndFor
      \State 返回 true
    \EndFunction
    \Function{GetMaxUsableDepth}{$x$}
      \State $g\gets\Call{PreferenceSequence}{x}$
      \For{$d=k$ downto $0$}
        \If{not \Call{IsSaturated}{$g[d], d$}}
          \State 返回 $d$
        \EndIf
      \EndFor
      \State 返回 $0$
    \EndFunction
  \end{algorithmic}
\end{algorithm}

\subsection{随机深度选择与 InsertKick}

插入时须记录 insert\_depth 并在各深度 bin 间均匀分布。插入前先取随机哈希 \(s(x)\in[0,k]\)，再令
\[
  \mathrm{depth}=\min(\Call{GetMaxUsableDepth}{x},\,s(x)).
\]
该式为随机深度与可用深度的 min 组合，对贪心上界形成软上限，使高深度 bin 更长时间保持未饱和。在目标 bin \(g_{\mathrm{depth}}(x)\) 中：若有空位，则写入 \(x\)、记录 x.insert\_depth = depth，并在 Local Query Router \(A[g_I(x)]\) 中登记探测下标 \(j\)；若无空位，则踢出 bin 内 insert\_depth 最小的元素 \(y\)，将 \(x\) 写入被踢槽位，再对 \(y\) 调用 ReInsert。整条踢出链至多触发 \(k\) 次，与 \(O(k)\) 交换上界一致\cite{bender2022otsh}。

\begin{algorithm}[htbp]
  \caption{InsertKick：k-kick 插入}
  \label{alg:insert-kick}
  \begin{algorithmic}[1]
    \State 输入 $x$（已为 $\pi(x)$）
    \State $\mathit{depth\_max}\gets \Call{GetMaxUsableDepth}{x}$
    \State $s_x\gets\Call{RandomHash}{x}\bmod(k+1)$
    \State $\mathit{depth}\gets\min(\mathit{depth\_max},\,s_x)$
    \State $g\gets\Call{PreferenceSequence}{x}$；$\mathit{target}\gets g[\mathit{depth}]$
    \If{$\mathit{target}$ 有空位（由 $M$ 判定）}
      \State 在任意空位写入 $x$；$x.insert\_depth\gets\mathit{depth}$
      \State 更新 $A,M,B$；返回成功
    \Else
      \State $y\gets\mathit{target}$ 中 $insert\_depth$ 最小元素；$\mathit{evict\_pos}\gets y$ 的位置
      \State 从 $\mathit{evict\_pos}$ 移除 $y$；将 $x$ 写入 $\mathit{evict\_pos}$
      \State $x.insert\_depth\gets\mathit{depth}$
      \State \Call{ReInsert}{$y$}
    \EndIf
  \end{algorithmic}
\end{algorithm}

\subsection{ReInsert：被踢元素向上放入}

插入发生踢出时，被踢元素须安置到合适深度；\textbf{ReInsert} 自 \(y.\)insert\_depth\(-1\) 向上逐层检查 Bin(\(y,d\))：若该 bin 在深度 \(d\) 下不饱和且有空位，或存在 insert\_depth 小于 \(d\) 的元素，则将 \(y\) 写入并可能继续踢出链式调用 ReInsert。深度降至 0 仍无法安置时 Cubby 已满，须新建 tail Cubby 并重试 InsertKick，与多层级 Cubby 设计衔接。MiniArray \(M\) 在此过程中的空位判定与定位均为常数时间。理论保证：在参数满足 Bender 等人\cite{bender2022otsh} 的设定时，极大概率至多 \(k\) 次 kick 后一定能安置被踢元素，如算法~\ref{alg:insert-kick} 所示。

\begin{algorithm}[htbp]
  \caption{ReInsert：被踢元素向上找空位}
  \label{alg:reinsert}
  \begin{algorithmic}[1]
    \Function{ReInsert}{$y$}
      \For{$d=y.insert\_depth-1$ downto $0$}
        \State $\mathit{bin}\gets\Call{Bin}{y,d}$
        \If{\Call{IsSaturated}{$\mathit{bin}, d$}}
          \State \textbf{continue}
        \EndIf
        \If{$\mathit{bin}$ 有空位（由 $M$ 判定）}
          \State 将 $y$ 写入空位；$y.insert\_depth\gets d$
          \State 更新 $A,M,B$；返回成功
        \ElsIf{$\mathit{bin}$ 中存在 $z.insert\_depth<d$}
          \State $z\gets\mathit{bin}$ 中 $insert\_depth$ 最小且满足 $z.insert\_depth<d$ 的元素
          \State 从 $\mathit{bin}$ 移除 $z$；将 $y$ 写入该槽位
          \State $y.insert\_depth\gets d$；更新 $A,M,B$
          \State \Call{ReInsert}{$z$}；返回
        \EndIf
      \EndFor
      \State 返回失败 \Comment{深度已降至 $0$ 仍无法安置，Cubby 已满}
      \State 新建 tail cubby；在新 tail 中对 $y$ 重试 \Call{InsertKick}{}
    \EndFunction
  \end{algorithmic}
\end{algorithm}

\subsection{复杂度与正确性}

单次 InsertKick 至多触发一条长度为 \(k\) 的踢出链：每次踢出仅调用一次 ReInsert，而被踢元素的 insert\_depth 严格下降，故链长受 \(k\) 约束，与 Bender 等人\cite{bender2022otsh} 的 \(O(k)\) 交换成本一致。

GetMaxUsableDepth 自 \(k\) 向 0 扫描，返回最深的不饱和 bin；随机深度 \(\mathrm{depth}=\min(\mathrm{depth\_max},s(x))\) 在理论上保证：若存在可写 bin，则插入不会触发深层踢出，同时使不同深度的 bin 负载趋于均衡\cite{mitzenmacher2001}。

bin 中的是否存在空位与定位空位都通过 MiniArray \(M\) 完成，避免对整个 Cubby 进行依次遍历 \(O(|I|)\) 时间复杂度的扫描；ReInsert 从 y.insert\_depth-1 向上逐层检查，每层仅访问 \(M\) 与局部 bin 状态，均摊常数时间。

元素每次写入或迁移，都需要同步更新 storage、insert\_depth、\(B\)、\(g_I(x)\) 中的 \(\pi(x)\mapsto j\) 映射对应的信息，以及 \(M\) 的空闲位图，这一同步规则是整个系统正常运行的前提，之后的操作才能正常运行。

\section{Local Query Router}

\subsection{设计目标}

偏好序列与探测位置序列 \(h_1(x),h_2(x),\ldots\) 在插入分析中完整定义了候选槽位顺序。然而 k-kick 插入会通过踢出动态挪动元素：元素实际位置一般不等于首选深度 \(g_k(x)\)，也不等于 \(h_1(x)\)。若查询仍沿 \(h_1,h_2,\ldots\) 线性扫描，时间复杂度将与探测序列长度同阶，无法达到 \(O(1)\)。因此必须为每个键保存精准映射
\[
  \pi(x)\longrightarrow j,
\]
其中 \(j\) 为探测序列中的下标，查询时直接访问下标 \(h_j(x)\) 对应的 storage 中的数据，再用 \(B\) 与 payload 三个数据恢复 \(\pi(x)\) 后与查询元素进行对比校验\cite{bender2022otsh,pagh2008retrieval}。

\subsection{映射规则}

Cubby 内定义局部首选槽位号（与 k-kick 偏好序列无关）：
\[
  g_I(x)=g^\ast(x)[1,\log|I|]\in[0,|I|-1].
\]
所有满足 \(g_I(x)=i\) 的键由第 \(i\) 号 Local Query Router \(A[i]\) 管理。查询时先计算 \(i=g_I(x)\)，再在 \(A[i]\) 中查询 \(\pi(x)\) 得到探测序列下标 \(j\)（若不存在则未命中）；随后访问 storage 中 \(h_j(x)\) 对应槽位，读取 payload 与 \(B\)，重建 \(g^\ast(x)\) 并校验是否命中。

\subsection{二叉前缀树存储与接口}

每个 \(A[i]\) 实现为一棵二叉前缀树；对外键为 \(\pi(x)\) 的比特串。接口为：
\begin{itemize}
  \item insert(key, j)：\(j\le\log K\)，仅需 \(O(\log\log n)\) 比特存储；如算法~\ref{alg:lqr-insert} 所示；
  \item query(key)：返回 \(j\) 或“不存在”；如算法~\ref{alg:lqr-query} 所示；
  \item delete(key)：删除映射。如算法~\ref{alg:lqr-delete} 所示。
\end{itemize}

树中不存完整键，仅存区分键的最短前缀；叶节点存 j\_value。

插入过程利用映射后的 \(\pi(x)\) 进行操作；下文以 \(x\) 表示 \(\pi(x)\) 的比特串。映射后的 \(x\) 的每一位（从高位到低位）依次决定在树中向左（0）还是向右（1）走。沿路径向下，若遇空子树则创建新节点；到达叶节点时标记为叶并存 j\_value。即使管理的键数较多时，树高仍受限于键的位数，故访问步数为 \(O(\log U)\)。

\begin{algorithm}[htbp]
  \caption{Local Query Router.insert}
  \label{alg:lqr-insert}
  \begin{algorithmic}[1]
    \Function{insert}{$\mathit{key}$, $j$}
      \State $\mathit{cur}\gets\mathit{root}$
      \For{$\mathit{key}$ 的每一位 $\mathit{bit}$}
        \If{$\mathit{bit}=0$}
          \If{$\mathit{cur.child0}$ 为空} 创建子节点 \EndIf
          \State $\mathit{cur}\gets\mathit{cur.child0}$
        \Else
          \If{$\mathit{cur.child1}$ 为空} 创建子节点 \EndIf
          \State $\mathit{cur}\gets\mathit{cur.child1}$
        \EndIf
      \EndFor
      \State $\mathit{cur.is\_leaf}\gets\mathit{true}$；$\mathit{cur.j\_value}\gets j$
    \EndFunction
  \end{algorithmic}
\end{algorithm}

\begin{algorithm}[htbp]
  \caption{Local Query Router.query}
  \label{alg:lqr-query}
  \begin{algorithmic}[1]
    \Function{query}{$\mathit{key}$}
      \State $\mathit{cur}\gets\mathit{root}$
      \For{$\mathit{key}$ 的每一位 $\mathit{bit}$}
        \If{$\mathit{bit}=0$}
          \If{$\mathit{cur.child0}$ 为空} 返回 NOT\_FOUND \EndIf
          \State $\mathit{cur}\gets\mathit{cur.child0}$
        \Else
          \If{$\mathit{cur.child1}$ 为空} 返回 NOT\_FOUND \EndIf
          \State $\mathit{cur}\gets\mathit{cur.child1}$
        \EndIf
      \EndFor
      \If{$\mathit{cur.is\_leaf}$} 返回 $\mathit{cur.j\_value}$
      \Else  返回 NOT\_FOUND
      \EndIf
    \EndFunction
  \end{algorithmic}
\end{algorithm}

\begin{algorithm}[htbp]
  \caption{Local Query Router.delete}
  \label{alg:lqr-delete}
  \begin{algorithmic}[1]
    \Function{delete}{$\mathit{key}$}
      \State 沿 $\mathit{key}$ 的比特路径找到叶节点
      \State 将叶标记为非叶（$\mathit{is\_leaf}\gets\mathit{false}$），清除 $\mathit{j\_value}$
    \EndFunction
  \end{algorithmic}
\end{algorithm}

\subsection{复杂度与 Patricia 压缩}

对键长 \(\log U=\Theta(\log n)\) 的比特串，query 与 insert 沿树向下至多 \(\log U\) 步；Patricia 风格的前缀压缩\cite{morrison1968,pagh2008retrieval,lehmann2025combined} 将仅有一位不同的键合并为同一压缩边，使实测 router 深度（router\_steps\_max）远小于键长上界。每条映射存 \(j\le\log K=O(\log\log n)\) 比特，与整体 \(O(\log^{(k)}n)\) 空间目标一致\cite{bender2022otsh,kuszmaul2024filters}。

动态维护。 k-kick 踢出链中，被搬移键须 delete 旧映射再 insert 新 \(j\)；删除操作清空叶 j\_value 但不破坏共享前缀路径上的内部节点，均摊开销与 \(A[i]\) 所管理的局部键数相关。在 \(K=\mathrm{polylog}\,n\) 且 Cubby 容量 \(|I|\le K\) 的设定下，单个 \(A[i]\) 的规模受控，router 更新与查询均摊为常数时间\cite{bender2022otsh,bender2024sosa}。

\section{商压缩与键恢复}

\subsection{设计动机与信息分工}

当前层次结构中，查询路由需用映射后二进制低位定位 Facility 与 Cubby 内数组；若槽位仍存完整键，则寻址路径上的信息上下文被重复占用。商压缩借助随机可逆置换 \(\pi\)，将键映射为 \(\pi(x)\)，并按下述原则拆分信息：\(\pi(x)\) 的低 \(\log N\) 位 \(g^\ast(x)=\pi(x)[1,\log N]\) 的各段位由全局路由与局部地址上下文（Facility 编号、Cubby 容量、实际槽位下标）及 MiniArray \(B\) 中的差分元数据共同承担；槽位 storage 仅保存无法由上述上下文直接推出的高位商，即 payload
\[
  \mathrm{payload}=\pi(x)[\log N+1,\log U].
\]
查询时 Local Query Router 给出探测位置序列 j , \(h_j(x)\) 定位到 storage 中的具体位置，然后从 storage\([h_j(x)]\)、\(B[h_j(x)]\) 与 Facility/Cubby 上下文联立恢复完整 \(\pi(x)\)，再经 \(\pi^{-1}\) 得到原始键。该分工使每个元素在槽位层面仅占用 \(\log U-\log N\) 比特的有效载荷，而将路由相关的 \(\log N\) 位分散到地址结构与变长元数据中，与整体 \(O(\log^{(k)}n)\) 空间目标一致\cite{bender2022otsh,pagh2008retrieval}。

\subsection{\texorpdfstring{\(\pi(x)\) 与 \(g^\ast(x)\) 的位段划分}{pi(x) 与 g*(x) 的位段划分}}

记 \(x[a,b]\) 为 \(x\) 自第 \(a\) 位至第 \(b\) 位（自低位向高位）组成的子串。对键 \(x\) 经置换得 \(\pi(x)\)，全局路由哈希为
\[
  g^\ast(x)=\pi(x)[1,\log N].
\]
在 \(g^\ast(x)\) 内部，自高向低继续划分为三段（见图~\ref{fig:gx-layout}，第二章已给出示意）：
\begin{itemize}
  \item Facility 编号段：\(r=g^\ast(x)[\log K+1,\log N]=\lfloor g^\ast(x)/K\rfloor\)，用于定位键所属 Facility；
  \item Cubby 中间段：\(g^\ast(x)[\log|I|+1,\log K]\)，用于在 Facility 内区分不同 Cubby 及局部桶范围；
  \item Local Query Router 段：\(g_I(x)=g^\ast(x)[1,\log|I|]\in[0,|I|-1]\)，与 k-kick 偏好序列无关，仅决定键归属 \(A[i]\) 的桶号 \(i\)。
\end{itemize}

k-kick 插入可能将元素 \(x\) 写入槽位 \(i\neq g_I(x)\) 的实际位置。此时 \(g_I(x)\) 无法由槽位下标 \(i\) 直接读出，须额外保存差分
\[
  \delta=g_I(x)-i,
\]
存入 MiniArray \(B[i]\)。

\subsection{MiniArray \texorpdfstring{\(B\)}{B} 中的 MetaEntry}

每个占用槽位 \(i\) 对应一条变长元数据 MetaEntry，经 MiniArray \(B\) 的 Access/Update 接口读写。逻辑字段如下：
\begin{itemize}
  \item delta：有符号整数 \(\delta=g_I(x)-i\)，编码为 zigzag varint，适应 k-kick 搬移后槽位与 router 桶号的不一致；
  \item mid\_bits：\(g^\ast(x)\) 在 Cubby 范围内的中间位，即 \(g_k=g^\ast(x)\bmod K\) 满足 \(g_k=mid\_bits\cdot|I|+g_I(x)\) 的商部分 \(\lfloor g_k/|I|\rfloor\)；位宽 \(\texttt{mid\_w}=\lceil\log_2((K-1)/|I|+1)\rceil\)，由 \((K,|I|)\) 派生；
\end{itemize}

编码顺序为：先写入 delta 的 varint，再写入 mid\_bits 的固定位宽字段；解码时按相同布局与 MetaCodecLayout 逆序读出。\(B[i]\) 的 bitlen 随条目实际长度变化，由 Facility 级 MiniArray 的 Rank/Select 机制定位，与第二章所述变长元数据维护方式一致。

\subsection{恢复：重建 \texorpdfstring{\(\pi(x)\)}{pi(x)}}

给定 Facility 编号 \(r\)、Cubby 容量 \(|I|\)、槽位下标 \(i\)、storage\([i]\) 中的 payload 及 \(B[i]\) 解码得到的 MetaEntry，按下式自底向上重建 \(g^\ast(x)\) 与 \(\pi(x)\)：
\begin{align}
  g_I(x) &= i+\delta, \\
  g_k &= \texttt{mid\_bits}\cdot|I|+g_I(x), \\
  g^\ast(x) &= r\cdot K+g_k, \\
  \pi(x) &= (\mathrm{payload}\ll\log N)\;\vert\; g^\ast(x).
\end{align}
各步均须校验 \(g_I(x)\in[0,|I|)\)、\(g_k<K\) 及 \(g^\ast(x)<N\)；任一不满足则恢复失败，视为未命中。原始键由 \(x=\pi^{-1}(\pi(x))\) 得到。

\begin{algorithm}[htbp]
  \caption{ReconstructGxPi：从槽位上下文恢复 \(\pi(x)\)}
  \label{alg:reconstruct-gx}
  \begin{algorithmic}[1]
    \Function{ReconstructGxPi}{$r$, $N$, $K$, $|I|$, slot, payload, meta, logN}
      \State $\mathit{g\_I}\gets \mathit{slot}+\mathit{meta.delta}$
      \If{$\mathit{g\_I}\notin[0,|I|)$} \State 返回失败 \EndIf
      \State $\mathit{g\_k}\gets \mathit{meta.mid\_bits}\cdot|I|+\mathit{g\_I}$
      \If{$\mathit{g\_k}\ge K$} \State 返回失败 \EndIf
      \State $\mathit{g\_star}\gets r\cdot K+\mathit{g\_k}$
      \If{$\mathit{g\_star}\ge N$} \State 返回失败 \EndIf
      \State 返回 $(\mathit{payload}\ll\mathit{logN})\;\vert\;\mathit{g\_star}$
    \EndFunction
  \end{algorithmic}
\end{algorithm}

\subsection{查询校验}

查询键 \(x\) 时，Router 给出探测下标 \(j\)，访问候选槽位 \(i=h_j(x)\)。只有当完整恢复的值与查询的值相同，才认为查询命中：

\begin{algorithm}[htbp]
  \caption{SlotPiMatches：槽位命中校验}
  \label{alg:slot-match}
  \begin{algorithmic}[1]
    \Function{SlotPiMatches}{cubby, slot, $r$, table, query\_gx}
      \If{cubby.slots[slot] 为空} \State 返回 false \EndIf
      \State meta $\gets$ \Call{DecodeMetaEntry}{cubby.B[slot]}
      \State gx $\gets$ \Call{ReconstructGxPi}{$r$, $N$, $K$, $|I|$, slot, payload, meta, logN}
      \If{gx 失败} \State 返回 false \EndIf
      \State 返回 gx $=$ query\_gx
    \EndFunction
  \end{algorithmic}
\end{algorithm}

Router 仅负责将查询映射到 \(O(I)\) 个候选位置；最终是否命中必须由 payload、\(B\) 与 \((r,N,K,|I|)\) 联立恢复的 \(\pi(x)\) 与查询侧计算的 \(\pi(x)\) 一致方可确认。该“Router 定位 + 商压缩校验”的两段式路径，是系统能够不存储完整键的前提\cite{bender2022otsh,pagh2008retrieval}。


\subsection{空间性质}

单个占用槽位的比特开销可分解为：（1）storage 中 \(\log U-\log N\) 比特 payload；（2）\(B[i]\) 中变长 MetaEntry，其中， \(\delta\) 采用 varint 编码，长度会随 \(|g_I(x)-i|\) 增大而增加，mid\_bits 占 \(\Theta(\log(K/|I|))\) 比特，router\_bits 与 insert\_bits 为常数宽度。Facility 编号 \(r\) 与 Cubby 容量 \(|I|\) 作为查询路径上的上下文信息，不重复写入每个槽位。整体上来看，商压缩将 \(\pi(x)\) 中与寻址重叠的低位从槽位载荷中抽离，使附加空间与 \(O\log^{(k)}n\) 级预算相容；MiniArray \(B\) 的变长编码则使元数据按条目实际所需位数记录，避免为所有槽位预留固定最大键长\cite{kurpicz2023pachash,lehmann2024shockhash}。

\section{插入、查询与删除}

插入。 计算 \(g^\ast(x)=\pi(x)[1,\log N]\) 与 Facility 编号 \(r\)；在对应 Facility 的 tail Cubby 内执行 InsertKick（算法~\ref{alg:insert-kick}），写入 storage、更新 \(M\) 与 \(B[i]\)，并在 \(A[g_I(x)]\) 中 insert(\(\pi(x),j\))。tail 满则挂入 1-tiered cubby 并新建 tail，必要时按算法~\ref{alg:tier-schedule} 触发合并\cite{bender2022otsh,mitzenmacher2001}。

查询。 计算 \(g^\ast(x)\) 定位 Facility 与 Cubby；令 \(i=g_I(x)\)，在 \(A[i]\) 中 query(\(\pi(x)\)) 得 \(j\)；访问 \(h_j(x)\)，用 payload、\(B\)、\(r\) 重建 \(g^\ast(x)\) 并与 \(\pi(x)\) 比较。该路径仅含常数次 MiniArray 与 Local Query Router 访问，查询时间为 \(O(1)\)\cite{bender2022otsh}。

删除。 定位元素所在 Cubby 与槽位，清空 storage，更新 \(A[g_I(x)]\).delete(\(\pi(x)\))、\(M\)与 \(B\)。若删除发生在非 tail Cubby，按算法~\ref{alg:tail-ops} 从 tail 随机抽取元素回填；回填写入同样走 k-kick 与 Local Query Router 更新路径，以恢复各 bin 的饱和与深度不变量。

表级扩缩容时，\(N\)、\(K\)、\(r_j\) 与商压缩上下文变化，须先恢复原始键再按新参数重插，而不能直接复制旧槽位的 payload 与 \(B\)\cite{bender2022otsh,bender2024sosa}。

\section{本章小结}

本章给出分层哈希表的主要结构与算法：Facility 级 MiniArray 以 B 树、Rank/Select 与 Split/Merge 维护变长元数据\cite{bayer1972,jacobson1989}；多层级 Cubby 与 tail 机制定义 tier 调度与删除回填不变量\cite{bender2024iceberg,pandey2023iceberght}；k-kick 经偏好序列、饱和判定、InsertKick 与 ReInsert 实现 \(O(k)\) 交换插入\cite{bender2022otsh,pagh2004cuckoo}；Local Query Router 经 \(A[i]\) 前缀树将查询降为 \(O(1)\)\cite{morrison1968,pagh2008retrieval}。四组件同步更新规则贯穿插入、踢出、删除回填与扩缩容重插；第四章报告各参数下的性能测量。
